\documentclass[11pt]{article}
\usepackage[T1]{fontenc}
\usepackage{lmodern}
\usepackage[margin=1in]{geometry}
\usepackage{amsmath,amssymb,amsthm,mathtools}
\usepackage{booktabs,longtable,array}
\usepackage{xcolor}
\usepackage[hypertexnames=false,colorlinks=true,linkcolor=blue!45!black,urlcolor=blue!45!black,citecolor=blue!45!black]{hyperref}
\usepackage{microtype,enumitem,fancyhdr,listings,needspace}
\hypersetup{pdftitle={iPRA without an introduction witness parameter},pdfauthor={Research report prepared for Vladimir Reshetnikov}}
\setlength{\emergencystretch}{2em}
\setlist{itemsep=3pt,topsep=5pt}
\newtheorem{theorem}{Theorem}
\newtheorem{lemma}[theorem]{Lemma}
\newtheorem{corollary}[theorem]{Corollary}
\newcommand{\Ncal}{\mathcal N}
\newcommand{\Fcal}{\mathcal F}
\newcommand{\Gcal}{\mathcal G}
\newcommand{\Jcal}{\mathcal J}
\newcommand{\ipra}{\mathrm{iPRA}}
\newcommand{\seq}[1]{\mathrel{\vdash_{#1}}}
\newcommand{\pred}{\operatorname{pred}}
\newcommand{\newest}{\nu_{\Pi}}
\pagestyle{fancy}
\fancyhf{}
\fancyhead[L]{\small iPRA without an introduction witness parameter}
\fancyhead[R]{\small 12 November 2026}
\fancyfoot[C]{\thepage}
\setlength{\headheight}{14pt}
\lstset{basicstyle=\small\ttfamily,breaklines=true,columns=fullflexible}
\title{\textbf{iPRA without an introduction witness parameter}\\[5pt]
\large A consequence-equivalent restriction of the native scalar calculus}
\author{Research report prepared for Vladimir Reshetnikov}
\date{12 November 2026}
\begin{document}
\maketitle

\begin{abstract}
An existential introduction normally supplies an arbitrary numerical term.
In the native scalar presentation of intuitionistic primitive recursive
arithmetic, that argument can be removed entirely: always use the newest
ambient numerical parameter, and use \(0\) when the ambient scope is empty.
All native consequences remain derivable, including consequences of arbitrary
open assumptions. The construction applies both to the eight-rule kernel
and to its seven-rule diagonal variant. Its ordinary realization is therefore
bi-interpretable with full first-order \(\ipra\), with whole numerical domains,
actual equality, and the same identity numerical round trips as the current
proposal.

The essential ingredients are a predecessor construction using equational
induction alone, a uniform naming formula obtained by numerical cases, and a
proof translation that regenerates premises after extending the parameter
scope. The last point avoids assuming that a completed proof with a
scope-dependent witness can simply be renamed or weakened.

The archive includes complete standalone Python checkers and matched controls.
In the supplied numeric syntax the seven-rule checker decreases from
\(932/1029/2770\) to \(924/1021/2753\) significant tokens / AST nodes /
normalized bytes. This is a modest, reproducible implementation reduction
supporting a new mathematical restriction. It is not a new global byte
record and is not a proof of absolute minimality. The new equivalence is
proved here in written mathematics; it has not been checked in Lean.
\end{abstract}

\section{Result and comparison target}

The comparison is pinned to commit
\[
\texttt{9f9545f5ffd9e0392d4a36639503aca12383450a}
\]
of \texttt{VladimirReshetnikov/ProveIt}, in
\texttt{Logic/PrimitiveRecursiveArithmetic/ConciseIPRA}.
The project already contains three-family clause calculi, computation
normalization, inferred witnesses at elimination, and Lean equivalence
theorems \cite{repo,diag,bridge}. These are established parts of the baseline.

The new restriction concerns the \emph{introduction} witness. Let
\(\Pi=(a_1,\ldots,a_m)\) be the ordered list of numerical parameters currently
in scope, and put
\[
 \newest=
 \begin{cases}
   a_m,&m>0,\\
   0,&m=0.
 \end{cases}
\]
A fresh eigenparameter is appended to this list. The context \(\Gamma\) is
a separate list of assumptions: its last assumption need not introduce the
newest parameter.

\begin{theorem}[No witness argument]\label{thm:main}
Let \(\Ncal\) be the current eight-rule native kernel. Let \(\Fcal\) have
exactly the same syntax and rules, except that DI by cases must use
\(w=\newest\), with no witness supplied by the proof.
For every finite scope \(\Pi\), native context \(\Gamma\), and native formula
\(A\) in that scope,
\[
 \Pi;\Gamma\seq{\Fcal} A
 \quad\Longleftrightarrow\quad
 \Pi;\Gamma\seq{\Ncal} A.
\]
Let \(\Gcal\) be the current seven-rule diagonal calculus, and let \(\Jcal\)
make the same deterministic-witness restriction in its single introduction.
Then
\[
 \Pi;\Gamma\seq{\Jcal} A
 \quad\Longleftrightarrow\quad
 \Pi;\Gamma\seq{\Gcal} A
 \quad\Longleftrightarrow\quad
 \Pi;\Gamma\seq{\Ncal} A.
\]
\end{theorem}

The formula translations in these equivalences are literally the identity.
No restriction to a protected public sublanguage, quotient equality,
normal-form terms, or selected class of contexts is involved.

\begin{corollary}\label{cor:bi}
The ordinary realizations of \(\Fcal\) and \(\Jcal\) are equiconsistent and
bi-interpretable with full first-order \(\ipra\). The interpretations use
the entire numerical domains and actual equality, and the numerical
composites are the identity graph \(x=y\).
\end{corollary}

This proves the weaker zero-or-variable introduction restriction posed
as Question 2 in the pinned README: the deterministic rule uses only those
witnesses. The improvement is that the certificate need not choose even
which variable to use.

\section{Complete presentation}

\subsection{Syntax and intended readings}

There is one numerical sort. Raw terms are
\[
 t ::= a\mid 0\mid S(t)\mid I(t,t;u.t),
\]
where only \(u\) is bound in the last child of \(I\). Formulas are
\[
 A ::= (s=t)\mid Q_y A\mid D_{x,y}A.
\]
The designated contradiction is \(0=S0\).
There are no implicit ordinary logical connectives or hidden structural
inference rules. In the ordinary realization the binders mean
\[
 Q_y B(y) \equiv B(0)\to\forall y\,B(y),
 \qquad
 D_{x,y}B(x,y)\equiv
 \exists x\bigl(B(x,0)\to\forall y\,B(x,y)\bigr).
\]
These readings establish soundness and compose the bi-interpretations.
They are not extra rules of the checker.

All substitutions below are capture-free. Bound letters can be renamed.
A term template \(h(v)\) has one distinguished free hole; its other
parameters are retained.

\Needspace{0.48\textheight}
\subsection{The eight-rule presentation \texorpdfstring{\(\Fcal\)}{F}}

The following is the complete primitive inventory.
Freshness in DE and either introduction means absence from the ambient
scope, context, conclusion, and all displayed outside parameters.

\begin{longtable}{@{}p{0.10\textwidth}p{0.84\textwidth}@{}}
\toprule
Rule & Premises and conclusion\\
\midrule
Hyp &
If \(A\in\Gamma\), infer \(\Gamma\vdash A\).\\[5pt]
\(I0\) &
Infer \(\Gamma\vdash I(0,a;u.h)=a\).\\[5pt]
\(IS\) &
Infer \(\Gamma\vdash I(Sc,a;u.h)=h[I(c,a;u.h)/u]\).\\[5pt]
Tr &
From \(\Gamma\vdash s=t\) and \(\Gamma\vdash h(s)=r\), infer
\(\Gamma\vdash r=h(t)\).\\[5pt]
QE &
From \(\Gamma\vdash Q_yB(y)\) and \(\Gamma\vdash B(0)\), infer
\(\Gamma\vdash B(t)\), for any term \(t\) in scope.\\[5pt]
DE &
From \(\Gamma\vdash D_{x,y}B\) and
\(\Pi,b;\Gamma,Q_yB(b,y)\vdash A\), infer
\(\Pi;\Gamma\vdash A\), where \(b\) is fresh.\\[5pt]
DI &
From
\(\Pi,z;\Gamma,B(\newest,0)\vdash B(\newest,Sz)\), infer
\(\Pi;\Gamma\vdash D_{x,y}B\), where \(z\) is fresh and
\(\newest\) is computed in the \emph{parent} scope \(\Pi\).\\[5pt]
QI &
For a literal equation \(E(y)\), from
\(\Pi,z;\Gamma,E(z)\vdash E(Sz)\), infer
\(\Pi;\Gamma\vdash Q_yE(y)\), where \(z\) is fresh.\\
\bottomrule
\end{longtable}

DI is supplied a matrix and a premise derivation, but no numerical term.
QE still permits arbitrary targets; the theorem does not restrict both
term arguments simultaneously.

Equality reflexivity is derived. For example \(I(0,t;u.u)=t\)
transported along itself with the identity template gives \(t=t\).
Symmetry, transitivity, and term congruence follow from Tr.

\subsection{The seven-rule presentation \texorpdfstring{\(\Jcal\)}{J}}

Retain Hyp, \(I0\), \(IS\), Tr, QE, and DE. Replace the two introductions
by one. Inspect the raw matrix \(B(x,y)\), before any witness is inserted:
\[
 \operatorname{diag}(B)
 \quad\Longleftrightarrow\quad
 B\text{ is literally }
 I(y,y;u.h)=I(y,y;v.k).
\]
The displayed occurrences of \(y\) must be the bound second coordinate,
not a free parameter with the same printed name. Set
\[
 \eta_B(z)=
 \begin{cases}
   z,&\operatorname{diag}(B),\\
   0,&\text{otherwise}.
 \end{cases}
\]
The sole introduction is
\[
 \frac{\Pi,z;\Gamma,B(\newest,\eta_B(z))
                 \vdash B(\newest,Sz)}
      {\Pi;\Gamma\vdash D_{x,y}B(x,y)}.
 \tag{DI\(^*\)}
\]
The matrix alone selects cases or progress; the scope alone selects the
witness. This is the existing diagonal introduction with its witness argument
deleted. It introduces no new guard, axiom, normalizer, or auxiliary syntax.

\section{Avoiding an invalid structural argument}\label{sec:portable}

A proof containing deterministic DI cannot in general be weakened
to a larger parameter scope merely by copying its tree. If a newer parameter
is added, the witness selected at an old DI node changes. Thus naming the
desired term and then silently weakening the old proof is insufficient.

We use proof-generating families. Fix a parameter scope and retained
assumptions. A \emph{portable proof family} for \(A\) is a procedure which,
for every finite extension by fresh parameters and additional assumptions,
constructs a derivation of the unchanged \(A\) in the larger scope.
Retained assumptions are realized by such families. Hypothesis lookup,
computation, and equational rules construct families directly.
No claim about moving a previously completed target proof is part of
this definition.

The completed constructions below take portable premises and return
portable results: in a larger scope they are run again, choosing fresh
local eigenparameters there. The final translation inducts on a
\emph{source} derivation, under an arbitrary substitution of its free
parameters and portable realizations of its hypotheses. It returns a
portable target proof family. Ordinary assumptions have such realizations
by Hyp.

A local case proof whose witness is a freshly chosen \(z\) is invoked
while \(z\) is newest. It is \emph{not} assumed portable with that fixed
\(z\). In any larger outer scope, the complete enclosing construction chooses
a new case parameter last. This distinction is used explicitly in naming.

\subsection{Cases and full formula transport}

The derived rules in this section transform the proof-generating families
just defined. Supplying a derived base means supplying its portable realization,
not applying an unproved cut rule to an arbitrary completed target proof.

In \(\Fcal\), conditional closure by cases is available for every formula:
\[
 \frac{\Gamma,B(0)\vdash B(Sz)}
      {\Gamma\vdash Q_yB(y)}.
 \tag{QC}\label{eq:qc}
\]
Introduce \(D_{r,y}B(y)\), with \(r\) unused, at the automatically selected
witness \(\newest\). Eliminate it at a fresh parameter and return the opened
\(Q\)-assumption. The case-premise constructor is run with \(z\) fresh
and newest in each requested outer scope.

In \(\Jcal\), the same construction works. If the matrix is diagonal,
\(B(\newest,0)\) is a theorem: both sides reduce by \(I0\) to \(0\).
Supply that theorem to the case premise and ignore the new progress
hypothesis. Otherwise DI\(^*\) is already in case mode. Thus QC is
available before general equation induction is recovered.

Portable full formula transport
\[
 s=t,\quad F(s)\quad\vdash\quad F(t)
 \tag{L}\label{eq:leibniz}
\]
follows by induction on \(F\). At equations it is equational reasoning.
At \(Q_yB(v,y)\), use QC: transfer the new assumption \(B(t,0)\) backwards,
apply the retained closure at \(Sz\), and transfer the successor row forwards.
At \(D_{x,y}B(v,x,y)\), open the premise at fresh \(b\), then introduce the
target at \(b\). At that point \(b\) is exactly the newest parameter.
Transfer its base backwards, apply the opened closure, and transfer its
successor row forwards. For a diagonal matrix in \(\Jcal\), its base is
provable, so the same case construction supplies progress while ignoring
the induction hypothesis.

This induction uses only transport for smaller formula templates.
Its input families are regenerated whenever a fresh parameter is added.

\subsection{Equation induction in the seven-rule version}

For a term \(s(y)\), define \(L_s(y)=I(y,y;u.s(y))\).
Unconditional cases, obtained from QC and QE, prove
\[
 C(y,s(0),L_s(y))=s(y),
 \qquad C(v,a,b):=I(v,a;u.b).
 \tag{Dec}
\]
At \(0\) this is the zero selector equation; at \(Sz\) it follows from
two successor computations.

Given the step \(s(u)=t(u)\vdash s(Su)=t(Su)\), retain \(s(0)=t(0)\).
Under \(L_s(u)=L_t(u)\), the decoders and the retained base give
\(s(u)=t(u)\). Apply the step and compute back to
\(L_s(Su)=L_t(Su)\).
This is a permitted diagonal introduction with an unused witness coordinate,
so its deterministic choice is irrelevant. Elimination gives the conditional
closure of the lift, whose base is reflexive. QE at \(Sz\) gives
\(s(Sz)=t(Sz)\) under the retained original base. QC now yields
\(Q_y(s(y)=t(y))\).

Thus \(\Jcal\), as well as \(\Fcal\), has portable induction for every
literal equation. The argument uses QC, computations and equational
reasoning, without a predecessor or an unrestricted witness.

\section{The noncircular naming construction}\label{sec:names}

\subsection{Predecessor before existential naming}

Appendix~\ref{sec:eq-predecessor} proves, from computation,
equational reasoning, and literal-equation induction alone, that a scalar
iterator term \(\pred\) satisfies
\[
 \pred(0)=0,\qquad \pred(Sn)=n.
 \tag{P}
\]
Hence \(Sa=Sb\) implies \(a=b\) before any arbitrary term has been named.

This dependency matters. The existing predecessor proof uses contextual
parity cases in its Lean organization. Simply importing coverage into the
restricted calculus would risk assuming the unrestricted introductions being
recovered. The appendix replaces those steps by unconditional equational
identities. Its central repair is
\[
 M(r+k)=I(k,M(r);u.c(u))
\]
together with Boolean interpolation
\[
 H(Z(Zy))=C(y,H(0),H(1)).
\]
Every induction in that construction has a literal equation as its invariant.

\subsection{A gate and a uniformly available witness}

For an equation \(E\equiv(s=t)\), write
\[
 G_v(E):=C(v,0,s)=C(v,0,t).
\]
Its zero row is provable and its successor rows are interderivable with
\(E\), by computation and equational reasoning. Define
\[
 H(y):=D_{x,v}\,G_v\bigl(C(y,0,Sx)=y\bigr).
 \tag{H}
\]

First prove \(H(0)\). Whatever witness \(\newest\) is selected,
the inner equation \(C(0,0,S\newest)=0\) computes to \(0=0\).
The successor-row proof ignores its base assumption.

Next prove \(H(Sz)\) in a scope whose newest parameter is \(z\).
The rule selects \(z\), and \(C(Sz,0,Sz)=Sz\) is again a computation
instance. The case parameter is appended after all surrounding parameters.

In \(\Jcal\) these matrices are always in case mode: their outer iterators
have seed \(0\), whereas the diagonal guard requires seed \(v\).
QC gives \(Q_yH(y)\), and QE with \(H(0)\) gives
\[
 \Gamma\vdash H(St)\quad\hbox{for every term }t.
 \tag{HN}
\]
This is an explicit QC construction followed by QE, not unrestricted
generalization on a completed proof. The complete \(H(St)\) construction
is portable because its case parameter is selected anew in each outer scope.

Open \(H(St)\) at fresh \(b\). The new assumption is
\[
 Q_vG_v\bigl(C(St,0,Sb)=St\bigr).
\]
Its zero row is a theorem. QE at \(1=S0\), followed by computations,
gives \(Sb=St\). Applying \(\pred\) and (P) gives
\[
 b=t. \tag{Alias}
\]
This extraction uses only the retained closure, QE and equational reasoning,
so the alias equation is portable while \(b\) is retained in scope.

One can also obtain a standalone naming formula
\[
 N_t:=D_{x,v}G_v(x=t).
\]
While \(b=t\) is available and \(b\) is newest, introduce \(N_t\), then
discharge the opening of \(H(St)\). This names \(t\) without supplying it
as an introduction witness.

The gate cannot be omitted: \(D_{x,v}(x=t)\) with \(v\) unused has the
vacuous reading \(\exists x((x=t)\to(x=t))\).

\section{Completeness by regeneration of premises}\label{sec:compiler}

Take a finite source derivation \(\pi\), a capture-free substitution
\(\sigma\) sending its ambient parameters to terms in a target scope,
and portable realizations of its context after \(\sigma\).
We construct a portable target proof of its conclusion after \(\sigma\).

Hypothesis, computation, term transport, and QE translate pointwise.
At DE, produce the translated major, append a fresh target parameter \(b\),
and invoke the induction hypothesis on the source branch, extending
\(\sigma\) to send its eigenparameter to \(b\). The new closure assumption
has a portable realization by Hyp. The conclusion is free of \(b\), so
target DE discharges it.

At an introduction with source witness \(w\), put \(t=\sigma(w)\).
Construct \(H(St)\) and open it at a new \(b\), obtaining \(b=t\).
The target introduction must and does select \(b\). Append its fresh
step parameter \(z\).

For a source case introduction, the target premise supplies \(B(b,0)\).
Full formula transport gives \(B(t,0)\). Use this as the realization of
the source branch's extra assumption when translating that \emph{strictly
smaller} source derivation. Its result \(B(t,Sz)\) is transported back to
\(B(b,Sz)\). The target introduction is legal, and the outer DE discharges
\(b\). All additional assumptions remain in the target context; there is
no operation substituting \(t\) into a completed target proof.

In the eight-rule translation, source QI uses target QI.
For the seven-rule translation, start with a derivation in the existing
diagonal calculus. Its guard is preserved and reflected by capture-free
substitution of free parameters: such substitution cannot manufacture a
bound occurrence of \(y\). Thus source and target modes agree.
In progress mode the same construction transfers \(B(b,z)\) to
\(B(t,z)\), applies the recursively translated step, and transfers back.
In case mode it is the preceding construction. The alias changes the
selected witness, not the raw bound-variable guard.

For portability, run this procedure in each requested extension.
No completed target derivation is copied through a new parameter scope.
Recursive source calls decrease derivation height; auxiliary full transport
decreases formula depth; naming and predecessor are fixed finite macros.
This gives a terminating proof translation.

Conversely every target primitive inference is an instance of its source
rule. The existing equivalence of the diagonal calculus with the native
kernel \cite{diag} completes Theorem~\ref{thm:main}.
Only \emph{after} this result do we inherit arbitrary substitution,
weakening and cut as admissible target rules. Direct copying of proof trees
was not their proof of admissibility.

\section{Why the interpretation requirement is met}

The project proves faithful equivalence of the native calculus with its
ordinary scalar realization \(T_{\mathrm{eq}}\), and a bi-interpretation of
\(T_{\mathrm{eq}}\) with \(\ipra\) \cite{repo,bridge}.
Both numerical domains are whole. An iterator translates to primitive
recursion with the counter ignored in its step; the reverse translation
compiles primitive recursive definitions into scalar iterators.
The project proves term and formula round trips, not merely
equiconsistency.

Our change gives an equal proof relation on the same terms, formulas and
contexts. The existing translations and all round-trip proofs therefore
remain valid. If \(\tau,\rho\) are the project's translations, its native
proofs of
\[
 A\dashv\vdash \rho(\tau(A))
\]
translate by Theorem~\ref{thm:main}; the other composite is unchanged in
\(\ipra\). The domain predicates and identity graph \(x=y\) are unchanged.
This proves Corollary~\ref{cor:bi}.

The artifact does not claim to re-run or extend the repository's Lean build.
It uses the pinned existing bi-interpretation theorem and gives a written
proof of the new identity comparison. Selected source files and exact URLs
are included to make that dependency auditable.

\section{Standalone checker and measured simplification}

The main executable \texttt{code/checker\_seven.py} implements \(\Jcal\);
\texttt{code/checker\_eight.py} implements \(\Fcal\).
Each is a complete Python module using only the standard \texttt{json}
library. It checks the input sequent's syntax and scope, then the proof,
and prints \texttt{ACCEPT} only on success.
Neither checker has a term normalizer or imports an arithmetic prover.

The syntax uses \(0\) for zero, negative integers \(-1,\ldots,-n\) for
ambient parameters, positive de Bruijn indices for bound variables,
\texttt{[0,t]} for successor, and \texttt{[c,a,h]} for iteration.
Formulas are \texttt{[s,t]}, \texttt{['q',B]}, and \texttt{[B]}.
Input fields are \texttt{variables}, \texttt{context}, \texttt{goal},
and \texttt{proof}. A deterministic DI instruction is
\texttt{['di', matrix, premise]}; its witness is \(-n\), which is zero
in empty scope.

Matched controls use the same representation, validators, other rules and
I/O. Their DI instruction additionally supplies and validates an arbitrary
term. The code reduction is directly tied to the mathematical restriction.

\begin{center}
\begin{tabular}{@{}lrrr@{}}
\toprule
Complete checker & Tokens & AST nodes & Normalized bytes\\
\midrule
Eight-rule control & 979 & 1086 & 2834\\
Eight-rule deterministic witness & 971 & 1078 & 2817\\
Seven-rule diagonal control & 932 & 1029 & 2770\\
Seven-rule deterministic witness & 924 & 1021 & 2753\\
\bottomrule
\end{tabular}
\end{center}

All four delivered programs are the complete measured normalized sources,
without omitted runtime helpers. The metric script uses CPython 3.12.14,
\texttt{ast.unparse}, significant lexical tokens, and a full AST walk.
Bytes include one final newline.

The reduction is modest: eight tokens, eight AST nodes, and seventeen bytes
in each matched pair. In the pinned tagged-syntax programs the same direct
change saves three bytes and one AST node but adds two tokens, because
constructing the selected tagged variable costs code.
Those separate comparisons are recorded in the evidence.

The compact coordinate frontier---256 tokens, 278 AST nodes, 516 bytes---
comes from different programs with different syntax contracts.
Our standalone raw Eq/Q/D checker does not beat it.
Seven rules also do not beat the existing three-family clause inventory.
The improvement proved here is elimination of the introduction witness
argument, with a matched complete-checker reduction and an equivalence on
unrestricted native sequents.

\section{Validation and exact limits}

Executable validation covers all rule families, nontrivial naming
scaffolding, malformed syntax, scoping, eigenparameter escape, and attempted
false equations. Deterministic certificates are translated by inserting
the selected witnesses and checked by matched controls and the pinned
repository checkers. Results are recorded in
\texttt{evidence/validation.json}.

The certificate for \(H(2)\) obtains an existential successor witness
through QC and QE even though no DI instruction carries an arbitrary term.
It checks the numerical-cases portion of naming, not the entire predecessor
construction or the metatheoretic equivalence.

Accordingly:
\begin{itemize}
\item The rule systems and new equivalence are mathematical claims with
written proofs.
\item Checker sizes and the included certificate results are reproducible
executable evidence.
\item No new Lean or Rocq formalization, exhaustive proof of Python
correctness, absolute minimum, or global byte record is claimed.
\end{itemize}

Primitive equational transport is retained. This proof cannot simply be
transferred to a three-family calculus whose transport itself is recovered
from packets introduced at arbitrary witnesses; that would introduce a
circular dependency.

\section{Further research questions}

\begin{enumerate}
\item Formalize the deterministic rules and the proof translation in Lean.
Source derivations with realized contexts are the required induction
interface; naively renaming target proof trees is insufficient.
\item Formalize the coverage-free predecessor as an equational-fragment
theorem, separately from full numerical coverage.
\item Can deterministic introduction coexist with the three-family Name or
inferred-opening calculus without circularly invoking unrestricted packet
introduction to recover transport?
\item Does this restriction improve the actual compact 256-token,
278-node or 516-byte records under matched contracts?
\item Quantify certificate growth: replacing a witness term by a naming
derivation makes the checker smaller but can make certificates larger.
\item Can scope determine witnesses while introduction modes also avoid a
matrix guard, with a direct proof of all native consequences?
\end{enumerate}

\appendix

% Research note for the parent agent. This is a section fragment, not a standalone document.
\section{Predecessor from equational induction alone}
\label{sec:eq-predecessor}

This section removes a possible circularity from the restriction of existential
witnesses. The predecessor term is constructed before any general term-naming
principle, and without numerical coverage, a separation axiom, existential
introduction, or a case split under a hypothesis such as $c=0$.

\paragraph{The equational interface.}
Use the term constructors $0$, $S$, and
$I(n,a;u.h)$, the two computation equations
\[
 I(0,a;u.h)=a,\qquad
 I(Sn,a;u.h)=h[I(n,a;u.h)/u],
\]
ordinary equational reasoning, and induction for a literal equation. Thus,
if $z$ is absent from the parameters and hypotheses $\Gamma$, then
\[
 \frac{\Gamma\vdash E(0)\qquad
       \Gamma,E(z)\vdash E(Sz)}{\Gamma\vdash E(t)}.
\]
The unconditional cases rule
\[
 \frac{\Gamma\vdash E(0)\qquad\Gamma\vdash E(Sz)}
      {\Gamma\vdash E(t)}
\]
is the special case that ignores the induction hypothesis. This rule does
\emph{not} replace or specialize hypotheses of $\Gamma$ containing an
arbitrary term $c$. Every use of cases below proves an unconditional identity
with the displayed case variable and the other parameters independent.

All definitions below are abbreviations for scalar iterator terms. Write
\[
 C(c,a,b)=I(c,a;u.b),\qquad
 x+y=I(y,x;u.Su),\qquad 1=S0.
\]
Then $C(0,a,b)=a$, $C(Sz,a,b)=b$,
$x+0=x$, and $x+Sy=S(x+y)$. Equation induction gives
$0+x=x$, $Sx+y=S(x+y)$, associativity and commutativity of addition.
For example the second identity is induction on $y$, commutativity is
induction on $y$ using that identity, and associativity is induction on its
last argument. Hence finite sums can be rearranged, and $Sx=x+1$.
The notation $2x$ and $3x$ means $x+x$ and $x+x+x$.

\paragraph{Zero test and parity.}
Put
\[
 Z(x)=C(x,1,0),\qquad P(n)=I(n,0;u.Z(u)).
\]
The equations $Z(0)=1$, $Z(Sx)=0$, and
\[
 Z(Z(Zx))=Zx
\]
follow by computation and unconditional cases on $x$. Consequently
\[
 P(0)=0,\quad P(Sn)=Z(Pn),\quad Z(Z(Pn))=Pn.
\]
The last equation is induction on $n$: the base computes, and the step
is the displayed triple-zero-test identity. It follows that
$P(n+2)=Pn$. Induction on $n$ now gives
\[
 P(2n)=0,\qquad P(3n)=Pn.
\]
For the first step use $2(Sn)=2n+2$; for the second use
$3(Sn)=3n+3$ and $P(m+3)=Z(Pm)$.

\paragraph{Boolean interpolation without conditional cases.}
For any term template $H(v)$, with arbitrary other parameters, unconditional
cases on $y$ prove
\begin{equation}
 H(Z(Zy))=C(y,H(0),H(1)).
 \tag{BI}\label{eq:boolean-interpolation}
\end{equation}
Indeed $Z(Z0)=0$ and $Z(Z(Sz))=1$. Thus if an equation $p=Z(Zp)$ is
already available, equational congruence turns \eqref{eq:boolean-interpolation}
into $H(p)=C(p,H(0),H(1))$. In particular an identity between two term
templates at $0$ and at $1$ holds at such a $p$. This uses a previously proved
unconditional identity; it never substitutes into a contextual hypothesis.

\paragraph{The three-cycle and the ramp.}
Define
\[
 c(u)=Z(u)+2P(u),\qquad M(n)=I(n,0;u.c(u)),\qquad
 J(k,a)=I(k,a;u.c(u)).
\]
Computation gives $c(0)=1$, $c(1)=2$, and $c(2)=0$.
The iterator equations give $M(0)=0$, $M(Sn)=c(Mn)$, and induction on $k$
gives the addition law
\begin{equation}
 M(r+k)=J(k,M(r)).\tag{M+}\label{eq:mod3-add}
\end{equation}
Its step is $M(r+Sk)=c(M(r+k))=c(J(k,Mr))=J(Sk,Mr)$.

Now put
\[
 R(n)=I(n,0;u.S(u+M(u))).
\]
We claim $M(Rn)=Pn$. For its induction step write $r=Rn$ and $p=Pn$,
and use the induction hypothesis $Mr=p$. By \eqref{eq:mod3-add},
\[
 M(R(Sn))=c(M(r+Mr))=c(J(Mr,Mr))=c(J(p,p)).
\]
For a Boolean $p$, the identity
\[
 c(J(p,p))=Zp
\]
follows from (BI): at $p=0$ both sides are $1$, and at $p=1$ both sides
are $0$, since $c(J(1,1))=c(c(1))=c(2)=0$.
The already proved equation $p=Z(Zp)$ applies. This proves the induction
step because $P(Sn)=Zp$. The base is immediate.
We therefore have the simpler ramp recurrence
\begin{equation}
 R(0)=0,\qquad R(Sn)=Rn+Pn+1.\tag{R}\label{eq:ramp-rec}
\end{equation}

Induction now gives
\begin{equation}
 R(2n)=3n,\qquad R(3n)=3Rn+Pn.\tag{R23}\label{eq:ramp-multiples}
\end{equation}
For the first step, $P(2n)=0$ and $P(2n+1)=1$, so two applications of
\eqref{eq:ramp-rec} increase $R(2n)$ by $3$.
For the second, put $r=Rn$, $p=Pn$, and $q=Zp$.
The three successive parities at $3n$, $3n+1$, $3n+2$ are $p,q,p$.
Using the induction hypothesis, the next left side is therefore
\[
 R(3n+3)=3r+p+3+p+q+p=3r+3p+q+3,
\]
which is $3R(Sn)+P(Sn)=3(r+p+1)+q$. No conditional parity split is used.

\paragraph{A factor map.}
Let $H(x)=C(Px,Rx,x)$. Then
\[
 H(2x)=3x,\qquad H(1)=1,\qquad H(3x)=3H(x).
\]
The first two follow immediately. For the third put $p=Px$, $r=Rx$ and
use \eqref{eq:ramp-multiples} to obtain
\[
 H(3x)=C(p,3r+p,3x)=C(p,3r,3x)=3C(p,r,x).
\]
The two identities used here,
\[
 C(v,a+v,b)=C(v,a,b),\qquad
 C(v,3a,3b)=3C(v,a,b),
\]
are unconditional identities proved by cases on $v$. In particular this
argument also avoids conditional parity cases.

\paragraph{Comparison.}
Define
\[
 F(y,a)=I(y,a;u.H(u)),\qquad W(x)=I(x,1;u.u+u),\qquad
 K(x,y)=P(F(y,W(x))).
\]
Three elementary inductions on $y$ give
\[
 H(F(y,a))=F(y,H(a)),\qquad
 F(y,3a)=3F(y,a),\qquad F(y,1)=1.
\]
Their steps use the iterator equation and, respectively, the preceding
induction hypothesis, $H(3u)=3H(u)$, and $H(1)=1$.
Since $W(0)=1$ and $W(Sx)=2W(x)$, we get
\begin{align*}
 K(0,y)&=1,\\
 K(Sx,0)&=0,\\
 K(Sx,Sy)
 &=P(H(F(y,2W(x))))\\
 &=P(F(y,H(2W(x))))\\
 &=P(F(y,3W(x)))
  =P(3F(y,W(x)))=K(x,y).
\end{align*}
Consequently $K(Sx,x)=0$ by induction on $x$.
The intended reading of $K(x,y)$ is the Boolean indicator of $x\leq y$,
but that semantic reading is not needed: only the three displayed equations
and their diagonal consequence are used below.

\paragraph{Saturation and predecessor.}
Put
\[
 B(k,u)=C(K(S(Su),k),u,Su),\qquad
 T(t,k)=I(t,0;u.B(k,u)).
\]
The comparison equations and the unconditional selector identity
$C(v,Sa,Sb)=S(C(v,a,b))$ give
\[
 B(Sk,Su)=S(B(k,u)),\qquad B(S(Sk),0)=1.
\]
Induction on $t$ proves the conjugacy
\begin{equation}
 T(St,S(Sk))=S(T(t,Sk)).\tag{T}\label{eq:saturation-conjugacy}
\end{equation}
At $t=0$ it says $B(S(Sk),0)=1$. For the step, unfold the left side and
rewrite its inner value by the induction hypothesis:
\[
 T(S(St),S(Sk))
 =B(S(Sk),T(St,S(Sk)))
 =B(S(Sk),S(T(t,Sk)))
 =S(T(St,Sk)).
\]
An induction on $n$, applying \eqref{eq:saturation-conjugacy} at $t=k=n$
in its step, now gives $T(n,Sn)=n$.

Finally define the scalar iterator term
\[
 \operatorname{pred}(n)=T(n,n).
\]
Its base equation is immediate, and its successor equation is
\begin{align*}
 \operatorname{pred}(Sn)
 &=B(Sn,T(n,Sn))=B(Sn,n)\\
 &=C(K(S(Sn),Sn),n,Sn)
  =C(K(Sn,n),n,Sn)=C(0,n,Sn)=n.
\end{align*}
Therefore
\[
 \boxed{\operatorname{pred}(0)=0,\qquad
        \operatorname{pred}(Sn)=n.}
\]
Applying this term context to $Sa=Sb$ proves $a=b$.

\paragraph{Dependency conclusion.}
Every induction above has a literal equation as its invariant. Every case
argument proves an unconditional equation and is itself an equation induction
that ignores its induction hypothesis. The construction therefore belongs to
the purely equational iterator fragment. It does not use the full numerical
coverage principle, contextual case splitting, formula induction, successor
injectivity, or a name for an arbitrary term. In a restricted native calculus
with the two iterator equations, equational transport, QE and direct
QI-progress, this entire fragment is available before DI is considered at all.
This is the needed noncircular predecessor lemma for the witness restriction.

\begin{thebibliography}{9}
\bibitem{repo} Vladimir Reshetnikov, \emph{Concise presentations of intuitionistic primitive recursive arithmetic}, README and research index, ProveIt, pinned commit \texttt{9f9545f5ffd9e0392d4a36639503aca12383450a}.\newline
\href{https://github.com/VladimirReshetnikov/ProveIt/tree/9f9545f5ffd9e0392d4a36639503aca12383450a/Logic/PrimitiveRecursiveArithmetic/ConciseIPRA}{Pinned project and README}
\bibitem{diag} ProveIt, \emph{NativeDiagonal.lean}, same commit. Defines the diagonal calculus and proves \texttt{gprf\_iff\_nprf}.\newline
\href{https://github.com/VladimirReshetnikov/ProveIt/blob/9f9545f5ffd9e0392d4a36639503aca12383450a/Logic/PrimitiveRecursiveArithmetic/ConciseIPRA/Lean/ConciseIPRA/NativeDiagonal.lean}{Pinned diagonal calculus and equivalence proof}
\bibitem{bridge} ProveIt, \emph{NativeBiInterpretation.lean}, same commit. Native/iPRA faithfulness, both formula round trips, and inconsistency equivalence.\newline
\href{https://github.com/VladimirReshetnikov/ProveIt/blob/9f9545f5ffd9e0392d4a36639503aca12383450a/Logic/PrimitiveRecursiveArithmetic/ConciseIPRA/Lean/ConciseIPRA/NativeBiInterpretation.lean}{Pinned native/iPRA bi-interpretation}
\end{thebibliography}
\end{document}

